% Generated by Listas/make_statements.py from pratica-lecture-05.md, item 22. Do not edit.
\textbf{Source.} Tadelis (2013), \emph{Game Theory: An Introduction}, Exercise 9.5 ``Campaign Spending'', pp.~188--189 {[}PDF 207--208{]}. Answer in the source: none (the book prints no solutions; the preface says about half are on Princeton UP's website, not checked). Reproduced verbatim.



\begin{quote}

\textbf{Campaign Spending:} Two political candidates are destined to play the following two-stage game. Assume throughout that there is no discounting (\(\delta = 1\)). First, they compete in the primaries of their respective parties. Each candidate \(i\) can spend \(s_i \ge 0\) resources on ads that reach out to voters, which in turn increases the probability that candidate \(i\) wins the race. Given a pair of spending choices \((s_1, s_2)\), the probability that candidate \(i\) wins is given by \(\frac{s_i}{s_1 + s_2}\). If neither spends any resources then each wins with probability \(\frac{1}{2}\). Each candidate values winning at a payoff of \(16 > 0\), and the cost of spending \(s_i\) is equal to \(s_i\). After each player observes the resources spent by the other and a winner in the primaries is selected, they can choose how to interact. Each can choose to be pleasant (\(P\) for player 1 and \(p\) for player 2) or nasty (\(N\) and \(n\), respectively). At this stage both players prefer that they be nice to each other rather than nasty, but if one player is nasty then the other prefers to be nasty too. The payoffs from this stage are given by the following matrix (where \(w > 0\)):



{\def\LTcaptype{none} % do not increment counter

\begin{longtable}[]{@{}ccc@{}}

\toprule\noalign{}

Player 1 ~Player 2 & \(p\) & \(n\) \\

\midrule\noalign{}

\endhead

\bottomrule\noalign{}

\endlastfoot

\(P\) & \(w, w\) & \(-1, 0\) \\

\(N\) & \(0, -1\) & \(0, 0\) \\

\end{longtable}

}



\begin{enumerate}

\def\labelenumi{\alph{enumi}.}

\item

  Find the unique Nash equilibrium of the first-stage game and the two pure-strategy Nash equilibria of the second-stage game.

\item

  What are the Pareto-optimal outcomes of each stage-game?

\item

  For which value of \(w\) can the players support the path of Pareto-optimal outcomes as a subgame-perfect equilibrium?

\item

  Assume that \(w = 1\). What is the ``best'' symmetric subgame-perfect equilibrium that the players can support?

\item

  What happens to the best symmetric subgame-perfect equilibrium that the players can support as \(w\) changes? In what way is this related to the role played by a discount factor?

\end{enumerate}

\end{quote}

